\documentclass[oneside]{scrbook}

% \usepackage[utf8]{inputenc}
\usepackage[T1]{fontenc}
\usepackage{version}
% \usepackage[latin1]{inputenc}
\usepackage[italian]{babel}
\usepackage{algpseudocode}
\usepackage{amsmath}
\usepackage{booktabs}
\usepackage{color}
\usepackage{amsfonts}
\usepackage{url}
\usepackage{algorithm}
\usepackage{chngcntr}
\usepackage[answerdelayed]{exercise}
\usepackage{graphicx}
\usepackage{caption}
\usepackage{listings}
\usepackage{tikz}
\usepackage{pgf}
\usepackage{colortbl}
\usepackage{listingsutf8}
\usepackage[pdftitle = {Dispense di esercizi}, pdfauthor = {Gianmarco Tedeschi e Alberto Marchesi}, pdfsubject = {Fondamenti di Informatica ($082746$ - A.A.~$2025-2026$)}, pdfcreator = {Gianmarco Tedeschi e Alberto Marchesi}, pdfproducer = {-}]{hyperref}
% \usepackage[automark,nouppercase]{scrpage2}
\usepackage{mdframed}
\usepackage{scrlayer-scrpage} % Necessary package for page styles

%\pagestyle{scrheadings}

%\usepackage{titlesec}
%\newcommand{\sectionbreak}{\clearpage}
\counterwithin{Exercise}{chapter}
\counterwithin{Answer}{chapter}

\algrenewcommand\algorithmicthen{\textbf{Allora}}
\algrenewcommand\algorithmicif{\textbf{Se}}
\algrenewcommand\algorithmicelse{\textbf{Altrimenti}}
\algrenewcommand\algorithmicend{\textbf{Chiudi}}
\algrenewcommand\algorithmicwhile{\textbf{Finch\'e}}
\algrenewcommand\algorithmicdo{\textbf{Esegui}}
\algrenewcommand{\algorithmiccomment}[1]{\hskip3em\texttt{//#1}}

\setlength{\parindent}{0in}
\setlength{\parskip}{10pt}

\newmdenv[%
        outerlinewidth=2,%
        roundcorner=10pt,%
        leftmargin=0,%
        rightmargin=0,%
        outerlinecolor=black,
        innertopmargin=\topskip,%
        splittopskip=\topskip,%
]{algo}


\ohead{\headmark}
\ihead{}
\chead{}
\ifoot{A.A.~$2025-2026$}
\cfoot{Fondamenti di Informatica (082746) - G. Tedeschi}
\ofoot{Pag.~\pagemark}
\setheadsepline{.4pt}
\setfootsepline{.4pt}

\definecolor{lbcolor}{rgb}{0.9,0.9,0.9}

\lstdefinestyle{cstyle}{%
  inputencoding=utf8/latin1,
  breaklines=true,
  breakatwhitespace=false,
  basicstyle=\small\ttfamily,
  language=c,
  frame=single,
  showstringspaces=false
}

\lstdefinestyle{cstylesmall}{%
  inputencoding=utf8/latin1,
  breaklines=true,
  breakatwhitespace=false,
  basicstyle=\scriptsize\ttfamily,
  language=c,
  frame=single,
  showstringspaces=false
}

\lstdefinestyle{console}{%
  inputencoding=utf8/latin1,
  breaklines=true,
  breakatwhitespace=false,
  basicstyle=\scriptsize\ttfamily,
  showstringspaces=false,
  backgroundcolor=\color{lbcolor}
}


\lstnewenvironment{clang}{\lstset{style=cstyle}}{}
\lstnewenvironment{clangsmall}{\lstset{style=cstylesmall}}{}

\addtocounter{chapter}{1}

\begin{document}

\chapter{Input/Output, operatori matematici e costrutto \texttt{if}}

Questa dispensa propone esercizi sulla scrittura di algoritmi, in linguaggio C, utili alla comprensione delle operazioni aritmetiche e logiche tra numeri e del costrutto condizionale \texttt{if}. Si introducono anche le due funzioni principali della libreria \texttt{stdio.h}: \texttt{scanf()} e \texttt{printf()}.

%\begin{figure}[ht!]
%		\centering
%		\includegraphics[width = 0.5\textwidth]{rtfm.jpg}
%		\caption{Un consiglio da seguire.}
%		\label{img:rtfm}
%\end{figure}

Nel caso in cui ci fossero delle domande riguardanti il funzionamento dei costrutti di base si pu\`o consultare come riferimento:
\begin{itemize}
		\item Il manuale del corso: ``Informatica: arte e mestiere'', D. Mandrioli, S. Ceri, L. Sbattella, P. Cremonesi, G. Cugola, McGraw-Hill Education;
		\item Un manuale di C on-line: \url{http://www.cplusplus.com/}.
\end{itemize}

\textbf{Struttura di un programma C.}

\lstinputlisting[style=cstyle]{src/basic.c}

La riga di codice \texttt{Corpo del programma} non viene considerata dal compilatore, in quanto circondata dai caratteri \texttt{/*} e \texttt{*/}, che indicano l'apertura e la chiusura di un commento, rispettivamente. Nel caso in cui si voglia commentare da un punto fino alla fine della riga, pu\`o essere anche usata l'espressione \texttt{//}. L'istruzione \texttt{getchar()} non fa parte della soluzione. Si tratta di un'istruzione bloccante per mettere l'elaboratore in attesa di un carattere da tastiera. Senza questa istruzione, o istruzioni equivalenti (e.g., la \texttt{system(``PAUSE'')}), l'esecuzione del programma termina immediatamente senza permettere all'utente di visualizzare l'output a video.

\section{Input e Output}

Le funzioni principali per la comunicazione tra programma e utente in linguaggio C
sono \texttt{scanf()} e \texttt{printf()}, entrambe fornite dalla libreria standard
\texttt{stdio.h}.

\begin{itemize}
  \item \textbf{\texttt{printf()}}: serve per stampare informazioni a video.
        La stringa di formato specifica come visualizzare i valori delle variabili.
        Ad esempio:
        \begin{clang}
          int a = 5;
          printf("La variabile a ha valore: %d\n", a);
        \end{clang}
        In questo caso, lo specificatore di formato \texttt{"\%d"} indica che
        verrà stampato un numero intero.

  \item \textbf{\texttt{scanf()}}: serve per leggere dati inseriti da tastiera.
        Anche qui si utilizza una stringa di formato per specificare il tipo di
        dato atteso, seguita dall’indirizzo di memoria della variabile in cui
        salvare il valore letto. Ad esempio:
        \begin{clang}
          int b;
          scanf("%d", &b);
        \end{clang}
        Lo specificatore \texttt{"\%d"} indica che ci si aspetta un intero e
        l’operatore \texttt{\&} fornisce a \texttt{scanf()} l’indirizzo della
        variabile in cui memorizzarlo.
\end{itemize}

I principali specificatori di formato sono:
\begin{itemize}
  \item \texttt{\%d} : intero decimale (\texttt{int}).
  \item \texttt{\%f} : numero in virgola mobile a precisione singola (\texttt{float}).
  \item \texttt{\%lf} : numero in virgola mobile a doppia precisione (\texttt{double}).
  \item \texttt{\%c} : singolo carattere (\texttt{char}).
  \item \texttt{\%s} : stringa di caratteri (array di \texttt{char} terminato da \texttt{'\textbackslash 0'}).
  \item \texttt{\%x} : intero in formato esadecimale.
  \item \texttt{\%o} : intero in formato ottale.
\end{itemize}

In generale, \texttt{printf()} corrisponde all’operazione di \emph{scrittura} su
output, mentre \texttt{scanf()} corrisponde a quella di \emph{lettura} da input.



\section{Operazioni matematiche}
Si assume che il lettore sia familiare con i tipi di dato numerici previsti dal C (e.g., \texttt{int}, \texttt{float}) e con i rispettivi di caratteri che di specifica del formato (e.g., \verb|"%d"|, \verb|"%f"|). Inoltre, le operazioni matematiche essenziali necessarie alla comprensione degli esercizi proposti in questa sezione sono:

\lstinputlisting[style=cstyle]{src/basic-math.c}

Prima di procedere oltre, il lettore deve aver compreso il significato di questo frammento di codice (e.g., provando a compilarlo e ad eseguirlo).

\section{Costrutto \texttt{if} e condizioni}

Il costrutto \texttt{if} codifica un ramo condizionale. Il linguaggio C segue la seguente sintassi:

\begin{clang}
  if (condizione)
    statement;
  [else statement;]
\end{clang}
dove le parentesi quadre indicano che la parte \texttt{\textbf{else} statement;} \`e opzionale. 
Come per tutti gli altri costrutti in C, se uno \texttt{statement} \`e una sola istruzione terminata da punto e virgola, non serve altro. Se invece uno \texttt{statement} \`e composto da pi\`u istruzioni terminate da punto e virgola, sar\`a necessario racchiuderlo tra parentesi graffe, ossia:

\begin{clang}
  if (condizione) {
    istruzione1;
    istruzione2;
    ...
  }
\end{clang}

La \texttt{condizione} \`e un'espressione booleana, ovvero un'istruzione che, quando valutata, risulta sempre in un valore pari a zero ($0$, falso) o uno ($1$, vero). Per comporre espressioni booleane si utilizzano operatori \emph{relazionali} e \emph{booleani}.

\paragraph{Attenzione:} si osservi che in C, l'operazione di assegnamento \texttt{a = 3} \`e diversa dall'operazione di confronto \texttt{a == 3}. La prima \`e sempre valutata come vera ($1$, uno), mentre la seconda, ovviamente, dipende dal valore memorizzato in \texttt{a}. Perci\`o:
\lstinputlisting[style=cstyle]{src/assegnamento_vs_confronto.c}


\subsection{Esercizi}

\begin{Exercise}[label={ex:difference}]
	Scrivere un programma che esegua la differenza di due numeri interi inseriti da tastiera.
\end{Exercise}

\begin{Answer}[ref={ex:difference}]
	Risoluzione del problema: dovremo innanzitutto leggere e memorizzare i due numeri. Una volta eseguita l'operazione di differenza dovremo stampare a video il risultato dell'operazione.
	
  \lstinputlisting[style=cstyle]{src/differenza.c}
\end{Answer}

\begin{Exercise}[label={ex:sconto}]
  Scrivere un programma che riceve in ingresso un prezzo (numero razionale) ed uno sconto (intero tra $0$ e $100$) da applicare, e restituisce il prezzo scontato e il risparmio ottenuto.
\end{Exercise}

\begin{Answer}[ref={ex:sconto}]
	Risoluzione del problema: leggiamo e memorizziamo la cifra e lo sconto. Dopodich\'e calcoliamo il prezzo scontato e, facendone la differenza con il prezzo iniziale, calcoliamo l'ammontare dello sconto. Infine stampiamo a schermo i due valori ottenuti.
	
  \lstinputlisting[style=cstyle]{src/sconto.c}
\end{Answer}
\begin{Exercise}[label={ex:inv_tempo}]
  Scrivere un programma che prende in ingresso un tempo espresso in ore, minuti e secondi e ne restituisce l'equivalente in secondi.
\end{Exercise}

\begin{Answer}[ref={ex:inv_tempo}]
	Risoluzione del problema: leggiamo i tre dati in ingresso e li memorizziamo. Per calcolare i secondi innanzitutto trasformiamo le ore in minuti (moltiplicando per $60$) e successivamente i minuti in ore.
	
  \lstinputlisting[style=cstyle]{src/inv_tempo.c}
\end{Answer}

\begin{Exercise}[label={ex:tempo}]
  Scrivere un programma che prende in ingresso un tempo espresso in secondi e ne restituisce l'equivalente nel formato ore, minuti, secondi.
\end{Exercise}

\begin{Answer}[ref={ex:tempo}]
	Risoluzione del problema: leggiamo e memorizziamo i secondi da convertire. Dividendo (per intero) per $60$ calcoliamo il numero di minuti e, calcolando il resto, otteniamo i secondi nel secondo formato. Ripetendo l'operazione con i minuti riusciamo a trovare il numero di ore e minuti che ci occorrono nel secondo formato. Stampiamo a schermo ore, minuti e secondi. 
  \lstinputlisting[style=cstyle]{src/tempo.c}
  
  Oppure in maniera più compatta:
  
  \lstinputlisting[style=cstyle]{src/tempo2.c}
\end{Answer}

\begin{Exercise}[label={ex:resto}]
  Scrivere un programma che prende in ingresso un prezzo in euro e restituisce il numero minimo di banconote utilizzando solo pezzi da $50$, $20$ e $5$ euro. Indicare anche la moneta rimanente.
\end{Exercise}

\begin{Answer}[ref={ex:resto}]
		Risoluzione del problema: leggiamo il prezzo da input, quindi dividiamo per $50$ il prezzo (per intero) per trovare il numero di banconote da $50$. Dividiamo poi ci\`o che rimane per $20$ (per intero) e poi ci\`o che rimane per $5$ per avere il numero di banconote da $20$ e $5$. Il resto in monete \`e ci\`o che rimane.

    \lstinputlisting[style=cstyle]{src/resto.c}
\end{Answer}
%\begin{Exercise}[label={ex:caratteri}]
  Scrivere un programma che prenda in ingresso un carattere (e.g., \texttt{a}, \texttt{b}, \texttt{c}) ed un numero intero positivo $n$. Il programma dovrà stampare a video il carattere inserito, formattato come carattere e come intero. Poi dovrà stampare a video il carattere $n$ posizioni successive a quello inserito. Ad esempio, se il carattere inserito è la lettera \texttt{a} ed $n = 3$, il programma dovrà stampare a video il carattere \texttt{d}.
\end{Exercise}

\begin{Answer}[ref={ex:caratteri}]
  \lstinputlisting[style=cstyle]{src/caratteri.c}
\end{Answer}

\begin{Exercise}[label={ex:distanza}]
  Scrivere un programma che calcoli la distanza tra due punti, $a$ e $b$, interi sulla retta $y = 0$.

  \begin{center}
    \begin{tikzpicture}
      \draw (0,0) -- (3,0);

      \fill[black] (1,0) circle (2pt) node [above] {a};
      \fill[black] (2,0) circle (2pt) node [above] {b};
    \end{tikzpicture}
  \end{center}

\end{Exercise}

\begin{Answer}[ref={ex:distanza}]
	Risoluzione del problema: chiediamo di inserire due punti all'utente. Calcolo la differenza delle due ascisse. Se \`e zero i punti coincidono, se \`e negativa dobbiamo restituire l'opposto, altrimenti restituiamo direttamente la differenza.

  \lstinputlisting[style=cstyle]{src/distanza.c}

\end{Answer}
\begin{Exercise}[label={ex:distanza_abs}]
  Scrivere un programma che calcoli la distanza tra due punti, $a$ e $b$, interi su un retta. Potete utilizzare la funzione \texttt{abs()} della libreria \texttt{stdlib.h}, che calcola il valore assoluto di un numero intero.

  \begin{clang}
    printf("abs(1-2) = %d", abs(1-2));
    //output: abs(1-2) = 1
  \end{clang}
\end{Exercise}

\begin{Answer}[ref={ex:distanza_abs}]
    \lstinputlisting[style=cstyle]{src/distanza-abs.c}
\end{Answer}
\begin{Exercise}[label={ex:basso_medio_alto}]
		Scrivere un programma che legga da input un numero intero e stampi su output:
		\begin{itemize}
		  \item la stringa \verb+basso+ se il numero \`e compreso tra $0$ e $3$;
		  \item la stringa \verb+MediO+ se il numero \`e compreso tra $4$ e $8$;
		  \item la stringa \verb+ALTO!+ se il numero \`e compreso tra $9$ e $10$;
		  \item la stringa \verb+Numero non valido+ altrimenti.
		\end{itemize}
		
\end{Exercise}

\begin{Answer}[ref={ex:basso_medio_alto}]
    \lstinputlisting[style=cstyle]{src/basso.c}
    Oppure nella versione pi\`u ottimizzata:
    \lstinputlisting[style=cstyle]{src/basso2.c}
\end{Answer}
\begin{Exercise}[label={ex:triangolo}]
		\Question Scrivere un programma che dati tre interi positivi valuti se essi possono essere i lati di un triangolo

		\Question Nel caso di risposta positiva al punto precedente si comunichi anche il tipo di triangolo (scaleno, isoscele, equilatero, rettangolo)
\end{Exercise}

\begin{Answer}[ref={ex:triangolo}]
		Risoluzione del problema: verifichiamo che la somma di ogni coppia sia minore del terzo lato e che la loro differenza sia maggiore del terzo.
		
    \lstinputlisting[style=cstyle]{src/triangolo.c}

		Risoluzione del problema: verifichiamo se \`e equilatero (ha tutti i lati uguali), poi se \`e scaleno (tutti diversi) altrimenti \`e isoscele. Per verificare se \`e rettangolo verifichiamo che valga il teorema di Pitagora tra il lato pi\`u lungo e gli altri due.
		
    \lstinputlisting[style=cstyle]{src/tipo_triangolo.c}

\end{Answer}
\begin{Exercise}[label={ex:bisestile}]
  Scrivere un programma che legga da tastiera un numero intero che rappresenta un anno (e.g., $2012$) e che determini poi se tale anno \`e bisestile o meno. Si pu\`o assumere che il numero intero letto da tastiera sia sempre valido (e.g., di $4$ cifre, positivo).

  Un anno \`e bisestile se \`e multiplo di $4$ ma non di $100$, oppure se \`e multiplo di $400$.
\end{Exercise}

\begin{Answer}[ref={ex:bisestile}]
    \lstinputlisting[style=cstyle]{src/bisestile.c}
    oppure
    \lstinputlisting[style=cstyle]{src/bisestile2.c}   
\end{Answer}
\begin{Exercise}[label={ex:logica}]
  	Assumendo che \texttt{c1} e \texttt{c2} siano due variabili di tipo char, che memorizzano rispettivamente i valori \verb|'e'| ed \verb|'m'| indicare, per ognuna delle espressioni logiche:
  
	  \begin{enumerate}
			  \item se l'espressione \`e vera o falsa (per i valori delle variabili sopra indicati);
			  \item se \`e sempre vera per qualsiasi valore che le due variabili possono assumere;
			  \item se \`e sempre falsa per qualsiasi valore che le due variabili possono assumere.
	  \end{enumerate}

  	Si fornisca una giustificazione per ogni risposta inserita nella tabella (risposte senza giustificazione potranno essere considerate nulle).

		\begin{enumerate}
			  \item \verb#((c1 != 'e') && (c2 == 'm')) || ((c1 != 'h') && (c2 == 'm'))#
			  \item \verb#(c1 < 'g') || !((c1 <= 'g') && (c1 != 'g'))#
			  \item \verb#(c1 <= 'm') || ((c2 > 'm') || !(c2 > c1))#
		\end{enumerate}
\end{Exercise}

\begin{Answer}[ref={ex:logica}]
\begin{enumerate}
	  \item \`E vera perch\'e \`e vero il secondo disgiunto (di cui \`e vero sia il primo congiunto, in quanto \texttt{c1} vale \verb|'e'|, che il secondo congiunto, in quanto \texttt{c2} vale \verb|'m'|). L'espressione, essendo vera, ovviamente non \`e sempre falsa. Inoltre non \`e sempre vera, per esempio sarebbe falsa se \texttt{c2} valesse \verb|'k'| o qualsiasi altro valore diverso da \verb|'m'|.
  	\item \`E sempre vera perch\'e, in base alla legge di De Morgan, \`e equivalente a
	  \begin{center}
	  		\verb#(c1 <'g') || (c1 > 'g') || (c1 == 'g')#
	  \end{center}
	  quindi la formula non \`e sempre falsa.
	  \item \`E sempre vera perch\'e equivalente alla formula
	    \begin{center}
	  		\verb#(c1 <= 'm') || (c2 > 'm') || (c2 <= c1)#
	  \end{center}
	  che \`e identicamente vera: infatti sarebbe falsa solo se fossero falsi tutti e tre i suoi disgiunti, cosa che non pu\`o essere, perch\'e se sono falsi i primi due (cio\`e se \verb#c1 > 'm'# e \verb#c2 <= 'm'#) allora \verb#c2 < c1#, quindi il terzo \`e vero. Quindi la formula non \`e sempre falsa.
\end{enumerate}
  
\end{Answer}

%to be checked
%\begin{Exercise}[label={ex:funzione_a_tratti}]
  Scrivere un programma che dato un numero reale $x$, calcoli la seguente funzione:
  \begin{equation*}
  	f(x) = \left\{ \begin{array}{lr}
  										x & x < 1\\
  										\exp{x} & 1 \leq x \leq 5\\
  										5 & x > 5
  			\end{array}\right.
  \end{equation*}
\end{Exercise}

\begin{Answer}[ref={ex:funzione_a_tratti}]
  \lstinputlisting[style=cstyle]{src/tratti.c}
\end{Answer}

\clearpage
\section*{Soluzioni}
\shipoutAnswer

\end{document}